\subsection{The formal full-state equation and its actual assumptions}
The retained architecture organizes the probe through an ordered equation
\begin{equation}
 (i\Gamma^\mu D_\mu-\Omega)\Psi=0,\qquad
 D_\mu=\partial_\mu+\mathcal B_\mu.
\end{equation}
Its ingredients include an internal connection, a closure projector,
a mismatch field and proposed triadic couplings. Specifying their
names does not determine their joint matrix elements on the generation
states. It also does not ensure the adjoint relations needed for a
real action or a selfadjoint time evolution.

On a chosen time chart with invertible $\Gamma^0$, the equation can
be rearranged algebraically as
\begin{equation}
 \partial_t\Psi=\mathcal G\Psi,
 \quad \mathcal G=-\mathcal B_0-(\Gamma^0)^{-1}
 \sum_{a=1}^3\Gamma^a(\partial_a+\mathcal B_a)
              -i(\Gamma^0)^{-1}\Omega.
\end{equation}
This expression retains the order of noncommuting factors. Turning it
into an evolution problem requires a domain, boundary or initial data,
regularity and an appropriate adjoint. In a finite modal comparison
these questions reduce to represented matrices; in a field theory they
cannot be removed by matrix notation.

The extraction condition given above tests an additional issue even
when a full evolution is already supplied. The reduced variables must
retain sufficient information to determine their own derivative.
The kernel of the extraction is the exact place where this condition
can fail. This is the dynamical counterpart of the electromagnetic
example, in which a Gram register omitted the orientation needed for
the transformed readout.

\subsection{A hidden-lift counterexample with equal normalization}
The dependence on discarded variables is not only a possibility for
large nonlinear systems. Consider a two-component complex amplitude
$f=x+ir$ evolving by $\dot f=-i\sigma_x f$, with $x,r$ real.
Then
\begin{equation}
 \dot x=\sigma_xr,\qquad\dot r=-\sigma_xx.
\end{equation}
Extract only $r$. Fix $0<b<1$ and compare the initial states
\begin{equation}
 r=(b,0)^{\mathsf T},\qquad
 x_\pm=(\pm\sqrt{1-b^2},0)^{\mathsf T}.
\end{equation}
Both full states have unit norm and the same extracted vector. Their
extracted derivatives are
\begin{equation}
 \dot r_\pm=(0,\mp\sqrt{1-b^2})^{\mathsf T}.
\end{equation}
The polar angle of $r$ therefore has opposite initial angular
velocities in the two cases. Equal extracted data and equal full
norm have not defined an autonomous angular law. This example is an
ordinary finite Hermitian evolution; it does not rely on an unstable
Hamiltonian or a noisy numerical reduction.

The linear criterion pinpoints the failure. If $T$ extracts $r$,
then $TG$ does not vanish on $\ker T$, because the discarded real
part drives $\dot r$. Adding a time-dependent extraction introduces
$\dot T$ rather than eliminating that condition. A moving frame may
be physically appropriate, but its derivative must be part of the
calculation.

\subsection{Exact elimination retains memory and initial-state data}
For a constant Hermitian modal Hamiltonian write
$H=S+iB$, where $S$ is real symmetric and $B$ real skew-symmetric.
The real equations are
\begin{equation}
 \dot x=Bx+Sr,\qquad\dot r=Br-Sx.
\end{equation}
Solving the first equation and substituting into the second gives
\begin{equation}
 \dot r(t)=Br(t)-Se^{Bt}x(0)
       -\int_0^tSe^{B(t-s)}Sr(s)\,ds.
\end{equation}
This is an exact reduced equation. It contains the discarded initial
state and a history kernel. Dropping either term defines a different
model unless an approximation and its range of validity are supplied.
In particular, a local damping coefficient cannot be inferred merely
from the fact that a component was projected out.

The same lesson applies to positive-root phase chambers. A projected
real or odd mode can remain on a kinematically selected slice while
its phase exits a positive-root chamber. A phase minimum of an energy
restricted to that slice is not enough to establish stability against
the omitted amplitude and complex-phase directions. The full action
and the resulting projected equation must determine which variations
are physically admitted.

\subsection{An action reduction that keeps the canonical partners}
A controlled finite alternative illustrates how additional inertia
can be derived without postulating a viscous law. Let
\begin{equation}
 L=\frac i2(f^*\dot f-\dot f^*f)-\mathcal E(f),\qquad
 e=(\cos(\theta/2),\sin(\theta/2))^{\mathsf T},
\end{equation}
and write
\begin{equation}
 f=(a+ib)e+cNe,\quad
 N=\begin{pmatrix}0&-1\\1&0\end{pmatrix},\quad b>0.
\end{equation}
The kinetic term is
$b\dot a-a\dot b-bc\dot\theta$. Removing $a$ and $c$ before
variation would remove the canonical partners of amplitude and
phase. Instead set $q=(b,\theta)$, $h=(a,c)$ and
$D_b=\operatorname{diag}(2,b)$. Up to a total derivative,
\begin{equation}
 L=-h^{\mathsf T}D_b\dot q-
       V(q)-j(q)^{\mathsf T}h-\tfrac12h^{\mathsf T}W(q)h.
\end{equation}
For positive definite $W$, the algebraic equation for $h$ is
\begin{equation}
 h=-W^{-1}(D_b\dot q+j),
\end{equation}
and completing the square gives
\begin{equation}
 L_{\rm eff}=\tfrac12(D_b\dot q+j)^{\mathsf T}
                    W^{-1}(D_b\dot q+j)-V.
\end{equation}
The induced kinetic metric is $D_bW^{-1}D_b$ and can depend on
the retained coordinates. Its derivatives belong in the equations
of motion. The hypotheses $W>0$ and $b>0$ are genuine conditions
for this regular reduction.

For the comparison $H=mI+\omega\sigma_y$, with $m>|\omega|$,
the result is
\begin{equation}
 L_{\rm eff}=\dot b^2/m+
       b^2(\dot\theta-2\omega)^2/(4m)-mb^2.
\end{equation}
The angular inertia is $b^2/(2m)$. Its equation
$\ddot\theta=-2(\dot b/b)(\dot\theta-2\omega)$ describes
reversible exchange with the amplitude. Contraction changes the sign
of the term resembling damping. This is why it cannot be interpreted
as positive viscosity without a separately represented dissipative
mechanism. These are modal comparison coordinates, not yet a native
fluid or a physical microscopic clock.

\subsection{Reality of the action and the meaning of a positive mismatch}
Consider ordinary Hermitian internal fields $A_\mu,H_c$, a fixed real
covector $n_\mu$, and
\begin{equation}
 \mathcal B_\mu=iA_\mu+\kappa H_cn_\mu.
\end{equation}
For constant $n$, write
\begin{align}
 D_\mu^AH_c&=\partial_\mu H_c+i[A_\mu,H_c],\\
 F^A_{\mu\nu}&=\partial_\mu A_\nu-\partial_\nu A_\mu
                                      +i[A_\mu,A_\nu],\\
 X_{\mu\nu}&=n_\nu D_\mu^AH_c-n_\mu D_\nu^AH_c.
\end{align}
The connection curvature is $F=iF^A+\kappa X$. Its literal quadratic
trace action and an adjoint-paired trace action differ:
\begin{align}
 \operatorname{Re}\left[-\tfrac14\operatorname{Tr}(FF)\right]
 &=\tfrac14\operatorname{Tr}(F^AF^A)
          -\tfrac{\kappa^2}{4}\operatorname{Tr}(XX),\\
 -\tfrac14\operatorname{Tr}(F^*F)
 &=-\tfrac14\operatorname{Tr}(F^AF^A)
          -\tfrac{\kappa^2}{4}\operatorname{Tr}(XX).
\end{align}
Spacetime indices are contracted in the declared real Lorentz metric.
The traced cross commutator in the second expression vanishes.
These formulas compare different actions, not two notations for the
same variational problem. The choice must be made consistently with
the probe law and its conserved currents.

Even positivity of $H_c$ supplies no general damping theorem. In
the constant-internal, flat-background comparison $H_c=hT_0$, with
$T_0$ Hermitian and $c_T=\operatorname{Tr}T_0^2>0$, the null
curvature term has the form
\begin{equation}
 L_{\rm HC}=\frac{\kappa^2c_T}{2}(n\cdot\partial h)^2,
 \qquad n^2=0.
\end{equation}
Its principal tensor is proportional to $n^\mu n^\nu$ and has
rank one for nonzero null $n$. A positive coefficient in one time
slice is not a four-dimensional Lorentzian wave operator. A physical
propagation or viscosity claim needs the rest of the represented
field system.

\subsection{Induced coframes and a direct same-frame obstruction}
Under a fixed Hermitian spin/internal factorization, let
$j^a=\Psi^*K^a\Psi$ and
$v_H^a=\Psi^*K^aH_c\Psi$. The proposed real coframe can then
have the identity
\begin{equation}
 E^a{}_{\mu}=\tfrac12\partial_\mu j^a+\kappa v_H^a n_\mu.
\end{equation}
If $H_c=0$ and $E$ has rank four, the functions $Y^a=j^a/2$
are local coordinates, and $g=\eta_{ab}dY^a dY^b$ is flat.
The inverse function theorem, rather than a curvature calculation,
establishes this restricted result. A variable adjoint, a geometric
spin connection or a different readout can change the hypotheses.

For the nonzero rank-one update $E=E_0+\kappa v n^{\mathsf T}$,
put $d=1+\kappa n^{\mathsf T}E_0^{-1}v$ and assume $d\ne0$.
The determinant lemma and inverse update give
\begin{equation}
 \det E=(\det E_0)d,\qquad
 n^{\mathsf T}g^{-1}n=
       \frac{n^{\mathsf T}g_0^{-1}n}{d^2},
 \quad g=E^{\mathsf T}\eta E.
\end{equation}
Thus a null covector remains null in this nonsingular update. This
is a geometric compatibility statement. A full solution must also
satisfy the probe, connection, closure and source equations.

There is a stronger consistency test if exactly the same probe,
adjoint and derivative are used both to reconstruct $E$ and to
define the coordinate gamma matrices in the ordered equation. Assume
\begin{equation}
 E^a{}_{\mu}=s_0\operatorname{Re}(\overline\Psi\Gamma^aD_\mu\Psi),
 \quad \Gamma^\mu=(E^{-1})^\mu{}_a\Gamma^a,\quad s_0\ne0,
\end{equation}
with real invertible $E$. Contraction of inverse and coframe yields
\begin{equation}
 4=s_0\operatorname{Re}(\overline\Psi\Gamma^\mu D_\mu\Psi).
\end{equation}
For $R_\Psi=(i\Gamma^\mu D_\mu-\Omega)\Psi$ and
$w=\overline\Psi\Omega\Psi$, it follows that
\begin{equation}
 \operatorname{Im}(\overline\Psi R_\Psi)=4/s_0-\operatorname{Im}w.
\end{equation}
If $w$ is real, the same-frame full-rank reconstruction cannot be
on shell. Reality of $w$ follows, for example, when the mass-like
operator has the required Hermiticity relative to the chosen probe
adjoint. This excludes a particular additional identification; the
Companion has not derived that identification as a native law. It
is not a universal inconsistency theorem for every geometric branch.

\subsection{Finite compatibility versus a field solution}
At a prepared finite jet, normal derivative equations can be written
as a linear system $Nx=r$ after lower data have been fixed. Let
$L_N$ have orthonormal columns spanning $\ker N^{\mathsf T}$.
Then the exact compatibility condition is
\begin{equation}
 L_N^{\mathsf T}r=0.
\end{equation}
For compatible $r$, the general solution is
$x=N^+r+k$, $k\in\ker N$. If a kernel decomposition is supplied,
its coordinates parametrize the remaining normal freedom. For
incompatible $r$, the least-squares value $N^+r$ leaves residual
$-L_NL_N^{\mathsf T}r$ and is not a field solution.

The distinction already appears in
$N=\begin{pmatrix}1&0\\0&1\\0&0\end{pmatrix}$, a three-by-two
matrix with first two rows the identity. A target $(a,b,c)$ is
compatible only when $c=0$, although a least-squares solver returns
$(a,b)$ for every $c$. A small normal residual or a solver success
flag cannot replace the explicit compatibility equation.

The source problem is more informative when higher coordinates split
into adjustable normal variables $w$ and prescribed tangential data $d$:
\begin{equation}
 h_0+Bq+A_fw+A_dd=0.
\end{equation}
Let $L_f$ span the left kernel of $A_f$ and define
\begin{equation}
 J_s=L_f^{\mathsf T}B,\qquad
 b_s(d)=-L_f^{\mathsf T}(h_0+A_dd).
\end{equation}
Eliminating $w$ leaves $J_sq=b_s(d)$. If $J_s$ has full column
rank, the source is unique for compatible data:
\begin{equation}
 (I-J_sJ_s^+)b_s(d)=0,\qquad q(d)=J_s^+b_s(d).
\end{equation}
The remaining higher lift is
$w=-A_f^+(h_0+A_dd+Bq)+k$, $k\in\ker A_f$.
This is a conditional finite source law. Changing which tangential
coordinates are prescribed changes the left-kernel test and can
remove a source constraint previously found by freezing all higher
variables. The derivative mask is therefore part of the mathematical
statement, not a numerical implementation detail.

\subsection{Constraint propagation requires open-set hypotheses}
The declared Cartan comparison uses an invertible real coframe, an
independent Lorentz connection, probe fields and multipliers enforcing
a constitutive coframe relation. Under compatible elimination of the
connection, coordinate and local Lorentz variations give residual
identities. Their source-free specialization requires both the
constitutive residual and the full probe residual to vanish on an
open region. A finite jet satisfying their point values does not
establish those hypotheses.

When the hypotheses do hold, and the spatially projected
trace-reversed coframe evolution equations are imposed, the
gravitational residual admits a symmetric hyperbolic subsidiary
system. With lapse $N_l>0$, shift $\beta$, spatial metric $h$ and
spatial covector $\xi$, its principal symbol can be written
\begin{equation}
 A(\xi)=-(\beta\cdot\xi)I_4+
 N_l\begin{pmatrix}0&(h^{-1}\xi)^{\mathsf T}\\
                         \xi&0_{3\times3}\end{pmatrix}.
\end{equation}
The positive matrix $W=\operatorname{diag}(1,h^{-1})$ symmetrizes
the symbol. Retaining the required lower-order terms and excluding
incoming constraint flux gives an energy estimate
\begin{equation}
 \mathcal E_g'(t)\leq C(t)\mathcal E_g(t).
\end{equation}
Zero initial constraints then remain zero under the stated regularity
conditions. This conditional propagation result does not propagate
the separate constitutive and probe equations from their initial
values. Those field equations are inputs to the subsidiary argument.

The geometric branch therefore has several distinct stages: an
algebraically meaningful coframe, a finite compatible derivative
system, a local field solution, and propagation of its constraints.
The programme has concrete tests for these stages, but a result at
one stage cannot be silently used as a result at the next. The same
discipline that separates positive mass readout from a physical mass
mechanism is needed when a finite source law is proposed as gravity.
